\documentclass[9pt,letterpaper]{article}
\usepackage[utf8]{inputenc}
\usepackage[T1]{fontenc}
\usepackage{geometry}
\geometry{top=0.55in,bottom=0.55in,left=0.68in,right=0.68in}
\usepackage{enumitem}
\setlist{itemsep=0pt, topsep=1pt, parsep=0pt}
\usepackage{hyperref}
\usepackage{setspace}
\usepackage{siunitx}
\sisetup{separate-uncertainty=true, range-phrase=--}
\DeclareSIUnit\angstrom{\text{\AA}}

\begin{document}
\fontsize{9}{10.5}\selectfont

\begin{flushleft}
\textbf{Myke Vital Sao Temgoua} \\
Department of Physics, Faculty of Science \\
University of Yaound\'{e} I \\
P.O. Box 812, Yaound\'{e}, Cameroon \\
\href{mailto:myke-vital.sao@facsciences-uy1.cm}{myke-vital.sao@facsciences-uy1.cm} \\[0.8em]
\today
\end{flushleft}

\begin{flushleft}
Prof.\ Kenneth M.\ Merz, Jr., Editor-in-Chief \\
\textit{Journal of Chemical Information and Modeling} \\
American Chemical Society
\end{flushleft}

\vspace{0.2em}

\noindent Dear Prof.\ Merz,

We resubmit the revised manuscript \textbf{``Target breadth and mutation resilience in African-natural-product-inspired antimalarial chemotypes: a computational analysis''} (original manuscript ID: ci-2026-00578g) as an Article in \textit{Journal of Chemical Information and Modeling}. We respectfully decline the transfer option and choose to remain with JCIM for the following reasons: (1)~JCIM is the premier venue for computational method development in cheminformatics and drug discovery, and this work directly addresses methodological questions in virtual screening validation, multi-target prioritization, and computational resistance profiling; (2)~the reviewers provided substantive, domain-expert feedback that significantly improved the manuscript, demonstrating that JCIM's editorial process is the right fit for this study; and (3)~JCIM's readership includes computational chemists working on resistance mechanisms and polypharmacology, the precise audience for this methodology.

We completed this revision in under 48~hours because we had already addressed the majority of the reviewers' concerns in parallel with the initial submission. The accompanying \textbf{point-by-point response document} demonstrates that all 18 reviewer comments (2 from Reviewer~1, 16 from Reviewer~2) have been resolved with explicit manuscript evidence, new analyses, or protocol clarifications. Major changes include: external DEKOIS~2.0 validation (two-arm comparison, MTX-retained vs MTX-stripped), a retrospective honest-negative control using five approved antimalarials, complete physicochemical profiling of all 17 candidates, PfCRT re-docking with corrected wild-type receptor and cavity-anchored grid, formal statistical multiplicity correction, and a staged Zenodo reproducibility package with sha256 verification.

This work addresses a methodological challenge in computational antimalarial discovery: evaluating target breadth and mutation resilience as distinct, complementary properties when prioritizing candidates from African natural-product-inspired chemical space, while maintaining clear boundaries between computational evidence and biological activity. A 17-member polypharmacology-oriented cohort was analyzed across four mechanistically diverse targets (PfDHFR, PfCRT, PfClpP, PfATP4) using target-specific structural anchors and a per-target resistance-resilience score (RRS). The workflow integrates three evidence layers: chemical-space expansion (\num{65856} hybrid-library molecules $\rightarrow$ \num{19913} prioritized candidates), target-anchored AutoDock Vina docking generating a complete $\num{17}\times\num{4}$ scoring matrix (\num{68}/\num{68} geometric quality pass), and per-target RRS computed separately for PfDHFR (N51I, C59R, S108N, I164L) and PfCRT (K76T, K76A) mutant panels. PfCRT was re-docked against a corrected 3D7-like wild-type receptor under a cavity-anchored grid; a pipeline-null control showed no detectable K76T/K76A signal for this panel. The cohort yielded seven A*, four B, and six C RRS profiles (range \qtyrange{73.2}{115.6}{\percent}). After multiplicity correction, PNS--RRS ($\rho=-\num{0.714}$), RRS--wild-type score ($\rho=\num{+0.691}$), and $N_{\mathrm{fav}}$--RRS ($\rho=\num{0.714}$) were significant; ACSI--RRS was not ($\rho=-\num{0.190}$).

\noindent\textbf{Key elements of the revision:}
\begin{enumerate}
\item \textbf{Protocol validation.} Three complementary checks: (a)~external DEKOIS~2.0 benchmark (PfDHFR ROC-AUC \num{0.45}, near-chance; controlled MTX-retained vs MTX-stripped comparison, \num{0.502} vs \num{0.563} with overlapping CIs, EF@1\% zero in both arms); (b)~MMV Malaria Box enrichment (ROC-AUC \numrange{0.924}{1.000}); (c)~redocking of reference ligands (four of five reproduced with RMSD \qty{<2.0}{\angstrom}; the methotrexate exception reflects NADPH-adjacent grid placement).
\item \textbf{Retrospective boundary check.} RRS profiles for five approved antimalarials (artemisinin, pyrimethamine, chloroquine, lumefantrine, cipargamin) against the same panels do not recover clinically documented resistance signatures (pyrimethamine $\approx$\qty{100}{\percent} on all four PfDHFR mutants; chloroquine $\approx$\qty{99}{\percent} on PfCRT K76T/K76A, both within the pipeline-null spread), reinforcing that docking-RRS values are not resistance-circumvention evidence.
\item \textbf{Physicochemical characterization.} All 17 candidates satisfy Lipinski/Veber criteria (\qty{100}{\percent} compliance; mean QED \num{0.703}); the ADMET summary is population-level and is not a candidate-specific safety assessment.
\item \textbf{Methodological transparency.} Complete grid specifications, within-target median favorability, pipeline-null RRS control, multi-seed sensitivity, and formal definitions of PNS and ACSI are provided. Centroid-based library reduction yields a \qty{99.3}{\percent} computational-cost reduction relative to exhaustive docking, with the caveat that this metric describes cost rather than active-compound retention.
\end{enumerate}

\noindent\textbf{Evidence boundaries.} The study does not claim experimental affinity, confirmed simultaneous target engagement, biological resistance circumvention, or measured IC\textsubscript{50}/EC\textsubscript{50} values. It produces (a)~four-target docking profiles from which polypharmacology hypotheses can be generated, (b)~per-target RRS values summarising mutation-panel docking ratios, and (c)~exploratory cross-metric associations without conflation of evidence layers.

\noindent\textbf{Reproducibility.} Complete records (source code, candidate lists, receptor and ligand files, grid configurations, raw poses, RRS protocols, 
machine-readable tables) are archived in the versioned Zenodo deposit (CC BY 4.0 license; 
\url{https://doi.org/10.5281/zenodo.22686176}). The Supporting Information contains the validation datasets and methodological specifications.

The centroid-based reduction strategy lowers the computational barrier for research groups in endemic regions (\qty{94}{\percent} of malaria cases) by orders of magnitude while retaining structural diversity, offering a reproducible template for hypothesis generation in resistance-challenged therapeutic areas. \textbf{We believe this revised manuscript now fully addresses all reviewer concerns and demonstrates the methodological rigor expected by JCIM's readership.} The comprehensive validation (DEKOIS, MMV, redocking, retrospective controls), transparent evidence boundaries, and complete reproducibility package represent the level of thoroughness appropriate for a methods-oriented cheminformatics journal. This manuscript is original, is not under consideration elsewhere, and has been approved by all authors. We declare no competing financial interests.

We thank the reviewers for their constructive feedback, which substantially improved the manuscript's clarity, validation rigor, and methodological transparency. We are grateful for the opportunity to resubmit to JCIM and look forward to your decision.

\vspace{0.3em}

\noindent Sincerely,

\noindent \textbf{Myke Vital Sao Temgoua} (on behalf of all co-authors)

\end{document}
